\chapter{Allgemeines}

\clearpage
\section{Pat.-Handling, Biomechanik}

\clearpage
\section{Medikamente allgemein}


\chapter{Teambuilding}

%\clearpage
%\section{Risiko- und Fehlermanagement}

\clearpage
\section{Crew Resource Management}

\chapter{Spezielle Notfälle}

\clearpage
\section{Akutes Koronarsyndrom}
\subsection{Definitionen}
\subsection{STEMI}
\subsection{NSTEMI}
\subsection{Lyse und Katheter}

\clearpage
\section{Rhythmologie} 

\clearpage
\section{Chirurgische KH Bilder}

\clearpage
\section{Thermische Zwischenfälle}

\clearpage
\section{Neurologie}
\subsection{Stroke und Lyse}
\subsection{Epilepsie}

\clearpage
\section{HNO}

\clearpage
\section{Urologie}

\clearpage
\section{Augen}

\clearpage
\section{Gynäkologie}

\clearpage
\section{Geburtshilfe}

\clearpage
\section{Tauchmedizin, HBO, CO}

\clearpage
\section{Kinder: Akute Erkrankungen}

\clearpage
\section{Kinder: Chirurgische Erkrankungen}

\chapter{Trauma}

\clearpage
\section{Unfallchirurgie Stamm}

\clearpage
\section{Unfallchirurgie Extremitäten}

\clearpage
\section{Neurotrauma}

\clearpage
\section{Kinder: Trauma}

\chapter{Anästhesie und Atemweg}

\clearpage
\section{Grundlagen der Beatmung}

\TextSlide{Unterschiede Atmung und Beatmung}
{%

}
{%
\begin{itemize}
  \item
\end{itemize}
}



\section{Zugänge, Drainagen} 
\subsection{Venenwege}

\subsubsection{Periphere Venenverweilkanülen}

\TextSlide{Periphere Venenverweilkanülen}
{%

}
{%
\begin{itemize}
  \item
\end{itemize}
}

\TextSlide{Periphere Venenverweilkanülen beim Kind}
{%

}
{%
\begin{itemize}
  \item
\end{itemize}
}

\subsubsection{Intraossäre Zugänge}

\TextSlide{Allgemein}
{%
Wenn ein intravenöser Zugang dringend erforderlich ist,
jedoch die Venensituation eine Anlage eines periphervenösen Katheters nicht 
zulässt,
ist ein intraossärer Zugang die Alternative der Wahl.
%
Es wurden zahlreiche kommerzielle Syteme zur Vereinfachung der Handhabung eines 
i.\,o.-Zuganges auf den Markt gebracht.
}
{%
\begin{itemize}
  \item
\end{itemize}
}

\TextSlide{Cook-Nadel}
{%

}
{%
\begin{itemize}
  \item
\end{itemize}
}

\TextSlide{EZio}
{%

}
{%
\begin{itemize}
  \item
\end{itemize}
}

\TextSlide{Bone Injection Gun}
{%

}
{%
\begin{itemize}
  \item
\end{itemize}
}

\subsection{Thoraxdrainagen}

\TextSlide{Allgemein}
{%
Eine Thoraxdrainage kann offen, im Rahmen einer (Mini-) Thorakotomie 
oder einer Thorakoskopie, 
aber auch geschlossen in Punktionstechnik über einen kleinen Hautschnitt 
und mittels eines Trokars durchgeführt werden. 
\cite{wiki:thoraxdrainge-20150127}




\subparagraph{Größe} Die Größe der Drainagen wird normalerweise in Charrière 
(CH) eingegeben, gängige Größen reichen bis zu 36 CH.

\subparagraph{Arten} 
Häufig verwendete Pleuradrainagen sind die Monaldi- und die Bülau-Drainage. 

\subparagraph{Bülau-Drainage}
Die 
Anlage der Bülau-Drainage erfolgt dabei in 
Höhe des 3. bis 5. Zwischenrippenraums (Intercostalraum, ICR) in der vorderen 
bis mittleren Axillarlinie, falls es sich um einen reinen Pneumothorax handelt, 
in Höhe der Schulterblattspitze in der mittleren bis hinteren Axillarlinie, 
falls auch Flüssigkeit (Hämatothorax, Pleuraerguss) abgesaugt werden soll. 
\cite{wiki:thoraxdrainge-20150127}

Triangle of saety: 

4-5. ICR, vordere Axillarlinie

\cite{youtube:ChestTubeInsertion}
\fullcite{youtube:ChestTubeInsertion}

\subparagraph{Monaldi-Drainage}
Der Punktionsort für die Monaldi-Drainage ist der 2. ICR 
in der Medioklavikularlinie.

Im Gegensatz zur Bülau-Drainage ist die Monaldi-Drainage meist  
kleinlumiger; sie kommt vor allem für die Behandlung eines Pneumothorax zum 
Einsatz. Für die Drainage der Luft 
ist das kleinere Lumen ausreichend und ermöglicht einen kleineren Hautschnitt 
mit weniger invasivem Eingriff.  \cite{wiki:thoraxdrainge-20150127}

\subparagraph{Spezielle Punktionssysteme}
TruClose
\cite{youtube:Tru-CloseThroracicVent}
\fullcite{youtube:Tru-CloseThroracicVent}

\subparagraph{Notfallmäßige Throaxpunktion}
Im Notfall, wenn in in ausreichender Zeit kein Drainagesystem verfügbar ist 
und ein lebensbedrohlicher Spannungspneumothorax 
besteht, 
kann dieser notafallmäßig durch Punktion des Thorax mit  
großlumigen Kanülen dekomprimiert 
und in einen offenen Pneumothorax umgewandelt werden.
%
Die akute Lebensbedrohung kann so kurzfristig abgewandt werden. 
%
Anschließend kann unter geordneten Bedingungen eine endgültige Drainage gesetzt 
werden. \cite{wiki:thoraxdrainge-20150127}



%Seite „Thoraxdrainage“. In: Wikipedia, Die freie Enzyklopädie. 
%Bearbeitungsstand: 3. Januar 2015, 14:27 UTC. URL: 
%http://de.wikipedia.org/w/index.php?title=Thoraxdrainage&oldid=137375584 
%(Abgerufen: 27. Januar 2015, 10:53 UTC) 
%
%@misc{ wiki:thoraxdrainge-20150127,
%   author = "Wikipedia",
%   title = "Thoraxdrainage --- Wikipedia{,} Die freie Enzyklopädie",
%   year = "2015",
%   url = 
%   "http://de.wikipedia.org/w/index.php?title=Thoraxdrainage&oldid=137375584",
%   note = "[Online; Stand 27. Januar 2015]"

\cite{wiki:thoraxdrainge-20150127}
}
{%
\begin{itemize}
  \item Offen:
  \begin{itemize}
    \item Minithorakotomie
    \item Thorakoskopie
  \end{itemize}
  \item Geschlossen: Punktionstechnik mit Trokar
\end{itemize}
}

\TextSlide{Anlage in Punktionstechnik}
{%
Das alternative Anlegen der Drainage durch Punktion mit 
einem Trokar birgt die Gefahr von Verletzungen des Lungengewebes und 
nachfolgenden Blutungen in sich.
}
{%
\begin{itemize}
  \item
\end{itemize}
}

\TextSlide{Anlage mittels Minithorakotomie}
{%
In der Regel wird die Thoraxdrainage mittels einer Inzision von 2 bis 3 cm 
(Minithorakotomie) angelegt. Nach der Inzision mit einem Skalpell und der 
Präparation mit einer Schere wird die zu drainierende Pleura mit dem Finger 
palpiert und gelöst. 
}
{%
\begin{itemize}
  \item
\end{itemize}
}

\section{Medikamente Anästhesie}


\section{Narkose}

\subsection{Rezepte}

\TextSlide{Narkose für ein Polytrauma ohne SHT}
{%
\begin{description}
  \item[Indikation] Polytrauma ohne SHT
  \item[Kontraindikation] SHT, 
  relevante kardiale Vorerkrankung
\end{description}
\begin{enumerate}
  \item \StepHeading{Präoxygenierung}
%  \begin{itemize}
%    \item
%  \end{itemize}
  \item \StepHeading{Einleitung}
  \begin{itemize}
    \item \E{Atropin} 0,5\Mg*
    \item Bezodiazepin: 
    \E{Diazepam} \Range{5}{10}\Mg*, 
    oder 
    \E{Midazolam} 5\Mg*
    \item \E{Esketamin} \Range{0,5}{1}\MgKg*
    \item Evt. Relaxierung: 
    \E{Lysthenon} 1\MgKg*
  \end{itemize}
  \item \StepHeading{Intubation}
%  \begin{itemize}
%    \item 
%  \end{itemize}
  \item \StepHeading{Aufrechterhaltung}
  \begin{itemize}
    \item \E{Esketamin} \Range{0,25}{0,5}\MgKg* alle 15\Min*\par
    oder 
    \E{Fentanyl} \Range{0,1}{0,2}\Mg* 
    (=\,\Range{2}{4}\Ml*) alle 20\Min*
    \item \E{Rocoronium} 0,5\MgKg* initial, 
    dann 0,15\MgKg* alle 20\Min*
  \end{itemize}
  \item \StepHeading{}Beatmung
  \begin{labeling}[:]{mmmmm}
    \item[AF]                  12\PerMin*
    \item[V\textsubscript{t}]  \Range{8}{10}\MlKg*
    \item[Fi\ce{O2}]           100\Pc
    \item[PEEP]                5
  \end{labeling}
\end{enumerate}
}
{%
\begin{itemize}
  \item 
\end{itemize}
}

\TextSlide{Narkose für ein SHT}
{%
\begin{description}
  \item[Indikation] SHT,
  relevante kardiale Vorerkrankung, 
  fortgeschrittenes Alter
  \item[Kontraindikation] 
\end{description}
\begin{enumerate}
  \item \StepHeading{Präoxygenierung}
%  \begin{itemize}
%    \item
%  \end{itemize}
  \item \StepHeading{Einleitung}
  \begin{itemize}
    \item \E{Fentanyl} \Range{0,1}{0,2}\Mg* (= \Range{1}{2}\Ml*)
    Benzodiazepin: \E{Midazolam} 5\Mg* \par
    oder
    \E{Diazepam} \Range{5}{10}\Mg*
    \item \E{Etomidat} 0,3\MgKg*
    \item Evtl. Relaxierung: \E{Lysthenon} 1\MgKg
  \end{itemize}
  \item \StepHeading{Intubation}
%  \begin{itemize}
%    \item
%  \end{itemize}
  \item \StepHeading{Aufrechterhaltung}
  \begin{itemize}
    \item 
    \E{Fentanyl} \Range{0,1}{0,2}\Mg* 
    (=\,\Range{2}{4}\Ml*) alle 20\Min*
    \item \E{Rocoronium} 0,5\MgKg* initial, 
    dann 0,15\MgKg* alle 20\Min*
  \end{itemize}
  \item \StepHeading{Beatmung}
  \begin{labeling}[:]{mmmmm}
    \item[AF]                  12\PerMin*
    \item[V\textsubscript{t}]  \Range{8}{10}\MlKg*
    \item[Fi\ce{O2}]           100\Pc
    \item[PEEP]                \Range{5}{10}
  \end{labeling}
\end{enumerate}
}
{%
\begin{itemize}
  \item 
\end{itemize}
}


\TextSlide{Thiopental}
{%
\begin{description}
  \item[Indikation] Polytrauma ohne SHT
  \item[Kontraindikation] SHT
\end{description}
\begin{enumerate}
  \item Präoxygenierung
  \begin{itemize}
    \item
  \end{itemize}
  \item Einleitung
  \begin{itemize}
    \item
  \end{itemize}
  \item Intubation
  \begin{itemize}
    \item
  \end{itemize}
  \item Aufrechterhaltung
  \begin{itemize}
    \item
  \end{itemize}
  \item Beatmung
  \begin{labeling}[:]{mmmmm}
    \item[AF]                  
    \item[V\textsubscript{t}]  
    \item[Fi\ce{O2}]           
    \item[PEEP]               
  \end{labeling}
\end{enumerate}
}
{%
\begin{itemize}
  \item 
\end{itemize}
}



\clearpage
\section{Geriatrischer Patient INT}

\clearpage
\section{Geriatrischer Patient UCH/ANÄ}

\clearpage
\section{Kinder Anästhesie}

\clearpage
\section{Status -A, -HT, DM, sonst.}

\chapter{Reanimation}

\clearpage
\section{CPR, Guidelines}

\clearpage
\section{Kinder: Reanimation}

\clearpage
\section{Alternative Reanimation }
(ZF 1995-2012)

\chapter{Großeinsatz und Gefahren}

\clearpage
\section{Führungsverfahren LSP}

\clearpage
\section{Stabsarbeit / LNA}

\clearpage
\section{Triage}

\clearpage
\section{Großschaden, MANV}

\clearpage
\section{CBRNE}

\clearpage
\section{Schadstoffe} 


\chapter{Rechtliche Grundlagen im Notarztdienst}

\clearpage
\section{Recht Grundlagen}

\clearpage
\section{Rechtssprechung}

\clearpage
\section{Recht und Urteile aus Arztsicht}
Einleitung

LSP I
